\section*{Chapitre \ref{ch:rs:half-life-decay-and-stability} --- Demi-vie, décroissance et stabilité}

\begin{solution}{exo:rs:half-life-decay-and-stability:1}
$-\ln 2/\ln 0.98 = 34.3$ jours~; un livre pondéré par les rangs échange environ $1 - 0.98 = 2~\%$ de lui-même par jour.
\end{solution}

\begin{solution}{exo:rs:half-life-decay-and-stability:2}
$1 - 0.37/0.85 = 56~\%$ (56.0~\% sur les moyennes non arrondies), proche de la baisse moyenne après publication de 58~\% observée par McLean et Pontiff.
\end{solution}

\begin{solution}{exo:rs:half-life-decay-and-stability:3}
$0.948(\sqrt{94} + 2/\sqrt{94}) = 9.39$ au premier résidu et $0.948(\sqrt{94} + 2 \times 94/\sqrt{94}) = 0.948 \times 3\sqrt{94} = 27.57$
au dernier.
\end{solution}

\begin{solution}{exo:rs:half-life-decay-and-stability:4}
Un signal d'âge $a$ dérive encore au jour $t + h$ si $a + h < 60$, ce qui, pour $a$ uniforme sur $0, \dots, 59$, a pour probabilité $(60 - h)/60$. L'IC
journalier à décalage $h$ est proportionnel à la part des signaux qui dérivent encore (chacun dérive au même rythme), il décroît donc linéairement et
s'annule en $h = 60$. Une loi exponentielle n'a pas de fin et une demi-vie constante~; forcée sur une droite qui s'arrête, elle annonce une demi-vie qui
ne décrit ni le début ni la fin.
\end{solution}

\begin{solution}{exo:rs:half-life-decay-and-stability:5}
$t = 0.04/(0.10/\sqrt{60}) = 3.10$. Selon la courbe de la \cref{fig:rs:half-life-decay-and-stability:posterior}, une division par deux mesurée a alors une probabilité a posteriori de
d'environ 0.86 (\texttt{rs\_decay.power(0.04, 0.10)} donne 0.858).
\end{solution}

\begin{solution}{exo:rs:half-life-decay-and-stability:6}
Chaque mois, le rendement réalisé du facteur déplace ensemble tous les titres chargés, dans le sens des rangs du prédicteur ou à l'encontre de ceux-ci~; la corrélation
de rang avec les rendements oscille alors avec le facteur, quelle que soit l'information propre à chaque titre que contient le prédicteur. Ces oscillations sont communes à toute la
section transversale~: elles ne se compensent pas entre les titres, seulement d'un mois à l'autre.
\end{solution}

\begin{solution}{exo:rs:half-life-decay-and-stability:7}
\texttt{rs\_decay.detection(0.5)}~: le sup-F détecte une baisse de 50~\% dans trois marchés sur dix et le CUSUM dans un seul, entre la baisse de 70~\% (sept et trois)
et l'absence de tout changement (un et un, les fausses alarmes). Sur un marché de dix ans avec quatre ans après la rupture, la division par deux de la dérive
d'un prédicteur est détectée moins souvent qu'une fois sur deux.
\end{solution}

\begin{solution}{exo:rs:half-life-decay-and-stability:8}
L'IC sur 24 mois du momentum est dominé par les rendements réalisés du facteur~: sur le marché synthétique, l'IC mensuel a un écart-type de 0.25
pour une moyenne de 0.02, si bien que deux ans ne disent presque rien (une division par deux n'est guère une preuve avec un $t$ de 0.72 sur cinq ans, et un tiers sur deux ans
encore moins). Le franchissement du CUSUM au dernier mois est exactement ce que le momentum synthétique a produit sans que rien ne change. Avant de tout arrêter~:
le taux de fausses alarmes du test sur l'ensemble des prédicteurs suivis, les indices de mécanisme, d'encombrement et de capacité (chapitre~28), et un échantillon plus long.
\end{solution}

\begin{solution}{pb:rs:half-life-decay-and-stability:1}
\begin{enumerate}
\item Retournement 0.037 et $-0.04$~; surprise 0.010 et 0.983~; momentum 0.008 et 0.995~; book-to-price 0.004 et 0.9995.
\item Son signal est le rendement du jour, non corrélé à celui de la veille ($\rho = -0.04$)~: les rangs changent chaque jour, et un livre qui les suit est
reconstruit chaque jour (rotation $1 - \rho = 104~\%$).
\item 20 jours~; la dérive injectée dure 60 jours.
\item Leurs IC journaliers (0.004 à 0.008) sont proches de leur bruit d'échantillonnage, et une demi-vie de plusieurs mois ne peut pas apparaître en 60 jours.
\item Momentum~: moyenne 0.023, écart-type 0.247. Surprise~: 0.040 et 0.041.
\item De $-0.035$ à 0.071.
\item Il se charge sur le facteur momentum, dont le rendement mensuel déplace l'ensemble de la section transversale.
\item Les deux tests rejettent au seuil de 5~\%~: le CUSUM franchit la limite au dernier mois, et le sup-F a une p-valeur de permutation de 0.023, alors que rien n'a changé.
\item 0.032 avant et 0.005 après.
\item Au mois 53, statistique 12.5, p-valeur 0.013. Les anciennes annonces continuent de dériver jusqu'à 60 jours après la rupture, si bien que l'IC baisse sur trois
mois.
\item Rien~: sa trajectoire atteint 0.69 de la frontière.
\item Dans ce marché, aucun des deux tests ne le voit (p-valeur du sup-F de 0.93)~; sur dix marchés, le sup-F le détecte dans sept cas et le CUSUM dans trois.
\item SMB 97~\%, HML 55~\%, Mom 56~\%~; p-valeurs du sup-F de 0.18, 0.15 et 0.053.
\item $p_0 = 0.37$ et $p_1 = 0.50$.
\item \textbf{Résultat nommé.} Avec un a priori d'un demi, une division par deux mesurée sur cinq ans signifie une vraie décroissance avec une probabilité de 0.57 pour un prédicteur
au bruit du momentum synthétique ($t = 0.72$), de 0.72 avec le bruit prépublication du facteur momentum ($t = 1.93$), et de 0.999 pour un prédicteur aussi stable que la
surprise synthétique ($t = 7.6$).
\item Environ 3.6~: avec $p_1 \approx \frac12$, la probabilité a posteriori vaut $0.5/(0.5 + p_0)$, qui atteint 0.9 lorsque $p_0 = 0.056$, soit
$\Phi(-t/\sqrt5) = 0.056$ (question d'entretien~6), $t = 3.6$.
\item Le mécanisme (existe-t-il une raison pour que cela se soit arrêté~?), les mesures d'encombrement et de capacité, la corrélation du prédicteur avec les nouveaux
prédicteurs publiés, les coûts de transaction, et le comportement du même prédicteur sur d'autres marchés.
\item Maîtriser les fausses alarmes sur l'ensemble des prédicteurs (un seuil plus strict, ou un taux de fausses découvertes, chapitre~20), exiger une persistance (deux
alarmes consécutives), et agir sur la taille avant de tout arrêter~: réduire le poids lorsque la preuve est faible, supprimer lorsqu'elle est forte.
\item La sélection de l'échantillon d'origine (une découverte fausse ou surestimée régresse), les changements de structure de marché (décimalisation, liquidité,
coûts), et un risque qui était rémunéré et l'est moins. McLean et Pontiff séparent le premier du trading en comparant les périodes hors échantillon et
après publication.
\item Lorsque le prédicteur est assez précis pour qu'une baisse de cette ampleur soit rarement due au hasard.
\end{enumerate}
\end{solution}

\begin{solution}{iq:rs:half-life-decay-and-stability:1}
Le décalage auquel l'\omterm{def:rs:half-life-decay-and-stability:decay}{autocorrélation du signal} tombe à la moitié, soit $-\ln 2/\ln\rho$ pour une autocorrélation d'ordre un $\rho$. Un livre qui suit le
signal échange environ $1 - \rho$ de lui-même à chaque période, si bien que les demi-vies courtes signifient une forte rotation et des coûts~; la demi-vie
propre à l'information, tirée de la \omterm{def:rs:half-life-decay-and-stability:decay}{courbe de décroissance de l'IC}, fixe la durée de détention, et les deux n'ont pas à concorder.
\end{solution}

\begin{solution}{iq:rs:half-life-decay-and-stability:2}
Calculer son IC avec le rendement de chaque période future unique $t + h$ (et non les rendements cumulés, qui mêlent les horizons), moyenner sur de nombreuses dates, puis tracer
en fonction de $h$~; n'ajuster une forme qu'après avoir regardé la courbe, avec des intervalles bootstrap sur les dates.
\end{solution}

\begin{solution}{iq:rs:half-life-decay-and-stability:3}
Comparer l'année avec le bruit propre de l'IC~: sa statistique $t$ annuelle, la probabilité d'une telle année si rien n'a changé, et un test de rupture
dont la p-valeur tient compte des queues épaisses (permutation) et du nombre de signaux suivis. Chercher ensuite d'autres indices~: mécanisme, encombrement,
capacité. Une seule mauvaise année pour un signal bruité est une preuve faible.
\end{solution}

\begin{solution}{iq:rs:half-life-decay-and-stability:4}
Biais statistique dans la découverte initiale (elle régresse hors échantillon), investisseurs qui tradent sur la publication, et changements de structure de marché
ou de coûts sans lien avec la publication. McLean et Pontiff distinguent les deux premiers en opposant la période hors échantillon mais antérieure à la publication (baisse de 26~\%,
borne supérieure du data mining) à la période postérieure à la publication (58~\%)~; la datation de la baisse par rapport aux \omterm{def:rs:half-life-decay-and-stability:break}{ruptures structurelles} (décimalisation) teste le
troisième.
\end{solution}

\begin{solution}{iq:rs:half-life-decay-and-stability:5}
Le CUSUM cumule les résidus récursifs standardisés et rejette lorsque la trajectoire sort de frontières qui s'élargissent avec le temps~; il n'a besoin d'aucune date candidate
et il est peu coûteux, mais sur la suppression injectée dans ce chapitre il n'a pas franchi la limite. Le sup-F essaie toutes les dates, est puissant contre un seul décalage
de la moyenne et le localise~; sa valeur critique doit tenir compte de la recherche (Andrews) ou être calculée par permutation. Les deux sont des tests à une seule rupture de la moyenne et
rejettent tous deux les séries stables et bruitées à leur taux nominal.
\end{solution}

\begin{solution}{iq:rs:half-life-decay-and-stability:6}
Soient les deux moyennes $\bar x_1 = \mu + \sigma_m z_1$ et $\bar x_2 = \mu + \sigma_m z_2$, avec $\sigma_m = \sigma/\sqrt n$ et des
normales centrées réduites indépendantes. L'événement $\bar x_2 \le \bar x_1/2$ s'écrit $z_2 - z_1/2 \le -\mu/(2\sigma_m)$, et $z_2 - z_1/2$ est normale de variance $5/4$, si bien que la probabilité
vaut $\Phi(-t/\sqrt5)$ avec $t = \mu/\sigma_m$~: 0.37 pour $t = 0.72$, 0.19 pour $t = 1.93$, 0.0004 pour $t = 7.6$.
\end{solution}
